\chapter{Introdução} \label{cap:introducao}

A crescente necessidade de promover acessibilidade, inclusão e autonomia em ambientes institucionais tem impulsionado o desenvolvimento de soluções tecnológicas voltadas ao controle inteligente de acesso a espaços físicos. Em especial, ambientes destinados a públicos com necessidades específicas, como pessoas autistas, neurodivergentes e gestantes, exigem não apenas segurança, mas também processos de acesso que reduzam barreiras sociais, cognitivas e sensoriais.

Nesse contexto, a Faculdade de Direito da Universidade de São Paulo (USP), localizada no histórico prédio do Largo de São Francisco, passou a contar com uma sala de apoio à amamentação e regulação sensorial como parte de suas políticas de inclusão e pertencimento. Apesar de sua importância, o uso desse espaço ainda é limitado por dificuldades relacionadas à segurança e à ausência de um sistema eficiente de controle de acesso. Atualmente, o modelo baseado na retirada manual de chaves junto a funcionários pode gerar constrangimentos, inconsistências de acesso e dificuldades adicionais, especialmente para usuários neurodivergentes.

Diante desse cenário, o projeto CAUSP-LOCK propõe uma solução integrada de hardware e software para viabilizar o acesso seguro, autônomo e simplificado ao ambiente. O sistema utiliza uma fechadura eletrônica baseada em microcontrolador ESP32-CAM e câmera OV2640, capaz de ler e validar QR Codes de acesso gerados por uma aplicação dedicada. Esses códigos são únicos, temporários e protegidos por mecanismos criptográficos, garantindo autenticidade, integridade e irretratabilidade dos acessos, além de permitir a operação offline da tranca.

Além de atender aos requisitos de segurança da informação, o sistema foi projetado com foco em acessibilidade, buscando minimizar a sobrecarga cognitiva e reduzir interações humanas obrigatórias no processo de acesso. Dessa forma, o CAUSP-LOCK se apresenta como uma solução que integra tecnologia, segurança e inclusão, contribuindo para o uso mais eficiente, seguro e digno de espaços institucionais compartilhados.


%%%%%%%%%%%%%%%%%%%%%%%%%%
\section{Motivação}
\begin{itemize}
    \item Subutilização de uma sala inclusiva com alto potencial de impacto social
    \item Dificuldades e constrangimentos no acesso devido ao modelo manual de retirada de chaves
    \item Falta de padronização e critérios claros para concessão de acesso
    \item Vulnerabilidade do espaço a furtos e uso indevido
    \item Necessidade de aliar segurança, acessibilidade e autonomia por meio de tecnologia
\end{itemize}

%%%%%%%%%%%%%%%%%%%%%%%%%%
\section{Objetivo}
\begin{itemize}
    \item Promover o bem-estar de pessoas autistas, neurodivergentes e gestantes
    \item Garantir o uso efetivo da sala de amamentação e regulação sensorial
    \item Reduzir a subutilização do espaço
    \item Aumentar a segurança contra furtos e uso indevido
    \item Eliminar a necessidade de retirada manual de chaves com funcionários
    \item Reduzir constrangimentos no processo de acesso
    \item Padronizar e formalizar o controle de acesso
    \item Permitir acesso autônomo aos usuários autorizados
    \item Implementar autenticação via QR Codes
    \item Garantir segurança por meio de criptografia (autenticidade e integridade)
    \item Utilizar códigos de acesso temporários e de uso único
    \item Possibilitar gestão de acessos via website ou aplicativo
    \item Minimizar esforço e sobrecarga cognitiva dos usuários
    \item Reduzir interações humanas obrigatórias no acesso
    \item Assegurar o uso adequado e controlado do espaço
\end{itemize}

%%%%%%%%%%%%%%%%%%%%%%%%%%
\section{Justificativa}
\begin{itemize}
    \item Necessidade de garantir acesso digno, autônomo e inclusivo a um espaço essencial
    \item Importância de reduzir barreiras enfrentadas por usuários neurodivergentes e gestantes
    \item Demanda por maior segurança e controle no uso de ambientes institucionais
    \item Aplicação de tecnologias acessíveis e de baixo custo para resolver problemas reais
    \item Contribuição para políticas de inclusão e melhoria da experiência dos usuários
\end{itemize}

%%%%%%%%%%%%%%%%%%%%%%%%%%
\section{Organização do Trabalho}
Descrever sucintamente os capítulos e as demais partes do trabalho.

Lembre-se de colocar rótulos nos capítulos e seções (com label{} ) e depois utilizar o ref{}
Por exemplo, o Capítulo \ref{cap:conceitos} apresenta conceitos....